% A ragged line listing delivered files (the paths are unbreakable).
\providecommand{\dirfiles}[1]{\par\smallskip\noindent
  {\footnotesize\raggedright\emph{Files:} #1.\par}\smallskip}
\providecommand{\dirP}{\mathbb{P}}
\providecommand{\dirE}{\mathbb{E}}
\providecommand{\dirU}{\mathrm{U}}

\section{Consequences through published theorems: primes, primitive roots and prime gaps}\label{sec:dir}

This section collects arithmetic consequences of Dirichlet GRH for single primes and prime gaps: the
abelian Dedekind zeta functions and the substitution of the input theorem into published conditional
theorems (\cref{dir:sec:published}), primitive roots and Euclidean division in cyclic fields
(\cref{dir:sec:roots}), a prime plus a multiple of a prime cube or fourth power (\cref{dir:sec:cube}), primes
in geometric intervals and Pomerance's divisibility conjecture (\cref{dir:sec:geometric}), the least
prime-generating constants for five exponents (\cref{dir:sec:mills}) and Toroczkai's prime-gap graphs
(\cref{dir:sec:graphs}). The other arithmetic consequences follow in this part: small primes in residue
classes, with subset products, sign patterns and codes, in \cref{sec:prd}; the class-group censuses of
imaginary quadratic fields and orders, Euler's idoneal list among them, in \cref{sec:cls}; and CM
coefficient fields, reflected fields and uniform Galois bounds in \cref{sec:cmf}.

Throughout this section the input theorem \cref{src:thm:dgrh} is used as a statement
without hypotheses: the Riemann zeta function and every Dirichlet $L$-function
$L(s,\chi)$, of any modulus and primitive or not, have no zeros in $\Re s>1/2$.
The results are of two kinds. Some follow by substituting the input theorem for the
analytic hypothesis of a published conditional theorem; each of these is stated once,
with the published theorem it rests on. The others require an additional argument,
which is given in full. A result marked ``\stG'' depends on \cref{src:thm:dgrh};
a result marked ``\stU'' does not. Several results use only the zeta case of the input
theorem, the classical Riemann hypothesis; this is said where it applies.

Logarithms are natural. For coprime integers $a,q$ with $q\ge1$ we write $P(a,q)$ for
the least prime congruent to $a$ modulo $q$, and for an odd prime $p$ we write $n(p)$
for the least positive integer that is not a square modulo $p$. For a fundamental
discriminant $D$, $\chi_D$ is the primitive Kronecker character and $h(D)$ the class
number of $\Q(\sqrt D)$.

\subsection{Abelian fields and consequences through published theorems}
\label{dir:sec:published}
\sources{grh_transfer_analysis; grh_transfer_analysis_2026-11-25; grh_lean_transfer_analysis; grh-transfer; grh-lean-transfers}

For a finite abelian extension $K/\Q$, class field theory attaches a finite group
$X_K$ of Dirichlet characters, the characters of $\Gal(K/\Q)$ viewed through the
Kronecker--Weber theorem. For $\chi\in X_K$ let $\chi^*$ be the primitive character
inducing $\chi$.

\begin{proposition}[Abelian Dedekind zeta functions; \stG]\label{dir:prop:abelian}
For every finite abelian extension $K/\Q$, every zero of $\zeta_K(s)$ in
$0<\Re s<1$ has real part $1/2$.
\end{proposition}

\begin{proof}
The factorization
\begin{equation}\label{dir:eq:abelian}
 \zeta_K(s)=\prod_{\chi\in X_K}L(s,\chi^*)
\end{equation}
holds for $\Re s>1$ by comparison of Euler factors and extends by meromorphic
continuation \cite{Sutherland2021Kronecker}. At an unramified prime $p$, if Frobenius has
order $h$, its values under the characters run through all $h$th roots of unity equally
often, so the local factor is $(1-p^{-hs})^{-[K:\Q]/h}$. At a ramified prime the
characters trivial on inertia supply the local factors for the quotient by inertia,
which is why each character is taken at its primitive conductor. The trivial character
supplies $\zeta(s)$. Every factor is holomorphic in the open strip, since its only
possible pole is at $1$. A finite product of nonzero complex numbers is nonzero, so a
zero of the product in the strip is a zero of some factor, and \cref{src:thm:dgrh}
places it on the line $\Re s=1/2$.
\end{proof}

The primitive characters matter for the exact identity \eqref{dir:eq:abelian}: writing
all factors at a common imprimitive modulus removes finite Euler factors. Those factors
do not vanish in the open strip, but they change the identity. The statement concerns
extensions abelian over $\Q$; an extension abelian over another number field leads to
Hecke $L$-functions outside the input theorem (see \cref{sec:hec}). In a cyclic field
of prime degree $\ell$ and conductor $f$ the nontrivial factors are the functions
$L(s,(\chi^j)^*)$, $1\le j\le\ell-1$, for one character $\chi$ of order $\ell$.

\begin{proposition}[Real nonprincipal characters; \stG]\label{dir:prop:positive}
Let $\chi$ be a real-valued nonprincipal Dirichlet character. Then
\[
 L(\sigma,\chi)>0\quad(\sigma>1/2),\qquad L(1/2,\chi)\ge0 .
\]
\end{proposition}

\begin{proof}
For real $\sigma>1$ the absolutely convergent Euler product has positive factors, so
$L(\sigma,\chi)>0$. Because $\chi$ is real, the analytic continuation is real-valued on
the real axis, and because $\chi$ is nonprincipal it is entire. By \cref{src:thm:dgrh}
it is nonzero on the connected interval $(1/2,\infty)$. A continuous real-valued
function cannot change sign there without vanishing, so it is positive throughout, and
continuity at $1/2$ gives the endpoint inequality.
\end{proof}

This removes real zeros in $(1/2,1)$, including possible exceptional zeros near $1$.
It does not give strict positivity at $1/2$, and zero location alone says nothing about
simplicity or multiplicity of zeros on the critical line.

\begin{theorem}[Published implications; \stG]\label{dir:thm:published}
The following statements hold.
\begin{enumerate}[label=\textup{(\alph*)}]
\item\label{dir:it:nonresidue} For every prime $p\ge5$, $n(p)<(\log p)^2$
 \cite[Corollary~1.1]{LamzouriLiSoundararajan2015Conditional}.
\item\label{dir:it:leastprime} For $q>3$ and $\gcd(a,q)=1$,
 $P(a,q)\le\bigl(\varphi(q)\log q\bigr)^2$
 \cite[Corollary~1.2]{LamzouriLiSoundararajan2015Conditional}.
\item\label{dir:it:lindelof} For primitive $\chi$ of conductor $q$ and every $\eps>0$,
 $\abs{L(1/2+it,\chi)}\ll_\eps (q(2+\abs t))^{\eps}$ \cite{Chandee2009Explicit}.
\item\label{dir:it:mobius} For every $\eps>0$, $\sum_{n\le x}\mu(n)=O_\eps(x^{1/2+\eps})$
 \cite{Soundararajan2009Partial}.
\item\label{dir:it:short} For every $x\ge2$ there is a prime in
 $(x-\tfrac4\pi\sqrt x\log x,\,x]$ \cite[Theorem~2]{Dudek2015Riemann}.
\item\label{dir:it:lagarias} For every $n\ge1$,
 $\sum_{d\mid n}d\le H_n+e^{H_n}\log H_n$, where $H_n=\sum_{k\le n}1/k$, with equality
 only at $n=1$ \cite[Theorem~1.1]{Lagarias2002Elementary}.
\item\label{dir:it:dbn} The de Bruijn--Newman constant equals zero.
\end{enumerate}
Items \ref{dir:it:mobius}--\ref{dir:it:dbn} use only the Riemann hypothesis for $\zeta$.
\end{theorem}

\begin{proof}
In \ref{dir:it:nonresidue}--\ref{dir:it:lagarias} the cited theorem assumes exactly the
Riemann hypothesis for $\zeta$ or for Dirichlet $L$-functions, which
\cref{src:thm:dgrh} supplies; the constants in \ref{dir:it:nonresidue} and
\ref{dir:it:leastprime} are those of the corrected third arXiv version of
\cite{LamzouriLiSoundararajan2015Conditional}. For \ref{dir:it:dbn}, let $H_t$ be the
usual heat deformation of Riemann's $\Xi$-function. The de Bruijn--Newman constant
$\Lambda_{\mathrm{DBN}}$ is characterized by all zeros of $H_t$ being real exactly when
$t\ge\Lambda_{\mathrm{DBN}}$, so the Riemann hypothesis is equivalent to
$\Lambda_{\mathrm{DBN}}\le0$. Rodgers and Tao proved $\Lambda_{\mathrm{DBN}}\ge0$
\cite{RodgersTao2020Bruijn}.
\end{proof}

Item \ref{dir:it:nonresidue} contains Vinogradov's conjecture
$n(p)\ll_\eps p^{\eps}$ \cite{Tao2015Elliott}, with an explicit constant.

\begin{corollary}[Vinogradov's least quadratic nonresidue conjecture; \stG]
\label{dir:cor:vinogradov}
For every $\eps>0$ and every odd prime $p$,
\[
 n(p)\le C_\eps\,p^{\eps},\qquad
 C_\eps=\max\Bigl\{2\cdot3^{-\eps},\ \frac{4}{e^2\eps^2}\Bigr\}.
\]
\end{corollary}

\begin{proof}
For $u\ge0$ the function $u^2e^{-\eps u}$ is maximal at $u=2/\eps$, with value
$4/(e^2\eps^2)$. Hence $(\log p)^2\le 4p^{\eps}/(e^2\eps^2)$, and
Theorem~\ref{dir:thm:published}\ref{dir:it:nonresidue} gives the bound for $p\ge5$. For $p=3$
the exact value $n(3)=2$ is covered by the other term of $C_\eps$.
\end{proof}

The inequality of Theorem~\ref{dir:thm:published}\ref{dir:it:nonresidue} was also checked by
exact arithmetic for all $428$ primes $5\le p<3000$, using Euler's criterion and
rational lower bounds for $\log p$. The bound in \ref{dir:it:leastprime} is roughly
quadratic in $q$; it is not Linnik's conjectural nearly linear bound.
\dirfiles{\artifact{grh_transfer_analysis/scripts/nonresidue_checks.py},
\artifact{grh_transfer_analysis/data/nonresidue_checks.json}}

\begin{remark}[What substitution does not give]\label{dir:rem:limits}
Three computations mark the boundary of this subsection.
First, Theorem~\ref{dir:thm:published}\ref{dir:it:short} does not place a prime between
consecutive squares: at $x=(n+1)^2$ the guaranteed interval has length
$\tfrac4\pi(n+1)\log((n+1)^2)$, and its ratio to the available gap $2n+1$ is asymptotic
to $(4/\pi)\log n$. Secondly, the Dirichlet-GRH estimate
$\psi(x;q,a)-x/\varphi(q)=O(\sqrt x\log^2(qx))$, summed over $q\le Q$, gives only
$O(Q\sqrt x\log^2(xQ))$; it does not reach the moduli of size nearly $x$ in the
Elliott--Halberstam conjecture, and a bound for one prime sum does not give the
two-prime correlations in the twin prime and binary Goldbach problems. Thirdly,
Hooley's conditional proof of Artin's primitive-root conjecture
\cite{Hooley1967Artin,Lenstra1977Artin,Perucca2026Closed} uses the Dedekind zeta
functions of the Kummer fields $\Q(\zeta_k,a^{1/k})$ for squarefree $k$, which are
generally nonabelian over $\Q$. Already for $k=3$ and $a=2$ the field is the splitting
field of $X^3-2$, with group $S_3$, and its zeta function has a two-dimensional Artin
factor not covered by Proposition~\ref{dir:prop:abelian}; see \cref{sec:grp}. These remarks
identify missing steps in known routes. They do not assert that the conjectures are
independent of the input theorem.
\end{remark}

\begin{remark}[A Turing machine that does not halt]\label{dir:rem:riemann95}
Let $M$ be the machine \code{riemann95} tabulated in the Cryptids repository
\cite{DeepMindCryptids}: it has $95$ working states, two tape symbols and one two-way
infinite tape, starts on the blank tape and halts on a missing transition. The
correctness theorem in that repository states, without hypotheses, that $M$ halts
exactly when Mathlib's \code{RiemannHypothesis} is false. Hence $M$ runs forever. The
two developments pin Lean 4.35.0-rc4 and 4.35.0, and the Mathlib file defining the
proposition is byte-identical at both pins; the composition is recorded as a
source-level example, and a joint build was not performed.
\dirfiles{\artifact{grh-transfer/checks/Riemann95Transfer.lean}}
\end{remark}

\subsection{Primitive roots and Euclidean division}\label{dir:sec:roots}
\sources{grh-lean-transfers; grh-transfer; grh_transfer_analysis_2026-11-25; grh_transfers}

For a prime $p$ let $g(p)$ be its least positive primitive root and $\widehat g(p)$ its
least prime primitive root.

\begin{theorem}[Grosswald's conjecture; \stG]\label{dir:thm:grosswald}
For every prime $p>409$, $g(p)<\sqrt p-2$. For every prime $p>2791$,
$\widehat g(p)<\sqrt p-2$.
\end{theorem}

\begin{proof}
This is Theorem~1 of McGown, Trevi\~no and Trudgian
\cite{McGownTrevinoTrudgian2016Resolving}, whose hypothesis is the Riemann hypothesis
for $\zeta$ and for the Dirichlet $L$-functions modulo $p$. Their proof detects a
generator of the cyclic group $\F_p^\times$ by smoothed prime sums. For $p\nmid n$ the
indicator of a primitive root is
\[
 \frac{\varphi(p-1)}{p-1}\sum_{d\mid p-1}\frac{\mu(d)}{\varphi(d)}
 \sum_{\ord(\chi)=d}\chi(n),
\]
where $\chi$ runs through the Dirichlet characters modulo $p$; this follows by writing
$n=u^k$ for a generator $u$ and applying inclusion--exclusion to $\gcd(k,p-1)=1$.
Their Lemma~1 bounds the sums of $\Lambda(n)(1-n/x)$ and of $\Lambda(n)\chi(n)(1-n/x)$
over $n<x$; the first uses the Riemann hypothesis and the second the hypothesis for
each $L(s,\chi)$. A prime power generating $\F_p^\times$ has a prime base that also
generates it, which gives the statement for $\widehat g$. Every character modulo $p$
is covered by \cref{src:thm:dgrh}.
\end{proof}

The endpoints are sharp: $g(409)=21$ and $\widehat g(2791)=53$. An unconditional proof
of the first bound for $p>1.77\cdot10^{54}$ is due to Bagger \cite{Bagger2026Note}.

\begin{corollary}[Erd\H{o}s problem 985 for odd primes; \stG]\label{dir:cor:erdos985}
For every odd prime $p$ there is a prime $q<p$ that generates $\F_p^\times$.
\end{corollary}

\begin{proof}
For $p>2791$ take $q=\widehat g(p)$ and use Theorem~\ref{dir:thm:grosswald}. For the remaining
$405$ odd primes, enumerate $p$ and the distinct prime divisors $\ell$ of $p-1$ by trial
division, and test successive primes $q<p$ until
\[
 q^{(p-1)/\ell}\not\equiv1\pmod p\qquad\text{for every prime }\ell\mid p-1 .
\]
The order of $q$ divides $p-1$, and a proper divisor of $p-1$ divides $(p-1)/\ell$ for
some such $\ell$, so the test certifies a generator. The program
\artifact{grh-lean-transfers/checks/primitive_roots.py} runs this test and
independently enumerates all powers of each witness. Every odd prime in the range
passes; the largest least witness is $53$, attained only at $p=2791$.
\end{proof}

The statement for all primes has the elementary exception $p=2$.

A number field $K$ is \emph{norm-Euclidean} if for all $a,b\in\OO_K$ with $b\ne0$ there
is $c\in\OO_K$ with $\abs{\Nm_{K/\Q}(a-bc)}<\abs{\Nm_{K/\Q}(b)}$; equivalently, every
$\xi\in K$ admits $\gamma\in\OO_K$ with $\abs{\Nm_{K/\Q}(\xi-\gamma)}<1$.

\begin{theorem}[Norm-Euclidean cyclic fields of prime degree; \stG]
\label{dir:thm:euclid}
The norm-Euclidean cyclic number fields of degree $5$ are the three fields of
conductor $11$, $31$, $41$, with discriminants $11^4$, $31^4$, $41^4$. Those of degree
$7$ are the two fields of conductor $29$, $43$, with discriminants $29^6$, $43^6$.
There is no norm-Euclidean cyclic field of degree
$19$, $31$, $37$, $43$, $47$, $59$, $67$, $71$, $73$, $79$ or $97$. A norm-Euclidean
cyclic field of degree $11$ has conductor $23$, $67$ or $331$.
\end{theorem}

\begin{proof}
These are Theorems~1.3, 1.4 and 1.5 of Lezowski and McGown
\cite{LezowskiMcGown2017Euclidean}, proved there under GRH. Their Proposition~5.2 uses
GRH bounds for character nonresidues and an auxiliary integer to bound the conductor of
a norm-Euclidean field; Proposition~2.1 and Table~2.2 give $f<6\cdot10^{10}$ in degree
five and $f<4\cdot10^{10}$ in degree seven, Proposition~2.2 supplies the finite
exclusions, and Proposition~2.3 the Euclidean minima of the surviving fields. The
conductor estimates go back to McGown \cite{McGown2012Norm}, who specifies the
hypothesis as GRH for the Dirichlet functions $L(s,\chi)$ or for $\zeta_K(s)$. For a
cyclic field $K$ of prime degree $\ell$ with character group generated by $\chi$,
\[
 \zeta_K(s)=\zeta(s)\prod_{j=1}^{\ell-1}L(s,(\chi^j)^*),
\]
so both forms of the hypothesis are supplied by \cref{src:thm:dgrh} and
Proposition~\ref{dir:prop:abelian}. The finite norm and exclusion computations are those of the
published source.
\end{proof}

The theorem gives no positive norm-Euclidean certificate for the three candidates in
degree eleven. As a check of one exclusion, the degree-$97$ field of conductor
$f=59\,999\,999\,974\,303$ has first two inert primes $2$ and $3$, and the integer
$431$ satisfies the character and congruence conditions of Proposition~3.1 in the arXiv
version of \cite{LezowskiMcGown2017Euclidean}. The corresponding classification of
cyclic cubic fields has a separate unconditional proof
\cite{BaggerEtAl2025Determination}.
\dirfiles{\artifact{grh-transfer/checks/check_euclidean_witness.py},
\artifact{grh-transfer/checks/euclidean-witness.txt}}

\subsection{A prime plus a multiple of a prime cube or fourth power}
\label{dir:sec:cube}
\sources{grh-arithmetic-transfers; grh_cm_prime_applications}

Lee and O'Clarey asked whether every integer greater than $24$ is the sum of a prime
and a positive multiple of a prime square, and proved this under Dirichlet GRH
\cite[Theorem~1.3]{LeeOClarey2026Prime}; an unconditional proof is reported by Campbell
\cite{Campbell2026Prime}. Raising the exponent gives sharp thresholds of a different
size. The quantifier on the multiplier matters: requiring the second summand itself to
be a prime power is a different problem.

\begin{theorem}[Prime cubes; \stG]\label{dir:thm:cube}
Every integer $n>10500$ has a representation $n=p+kq^3$ with $p,q$ prime and $k\ge1$.
Exactly $62$ positive integers have no such representation:
\[
\begin{gathered}
1,2,3,4,5,6,7,8,9,12,14,16,17,20,22,24,28,36,48,52,54,60,62,\\
72,76,78,96,102,108,120,126,150,174,180,216,240,270,294,312,342,\\
420,468,486,528,552,720,750,780,828,870,924,966,1080,1236,1260,\\
1500,1590,1800,1860,2520,8190,10500 .
\end{gathered}
\]
\end{theorem}

\begin{proof}
Let $q$ be the least prime that does not divide $n$, and
$B_q=\prod_{\ell<q,\ \ell\ \mathrm{prime}}\ell$. Then $B_q\mid n$ and $n$ is a unit
modulo $q^3$. If $P(n,q^3)<n$, subtraction supplies the required positive multiple of
$q^3$. By Theorem~\ref{dir:thm:published}\ref{dir:it:leastprime} applied to the modulus $q^3$,
\[
 P(n,q^3)\le\bigl[\varphi(q^3)\log(q^3)\bigr]^2<9q^8 .
\]
For every prime $q\ge43$ we have $B_q>9q^8$. Exact multiplication verifies this for
$43\le q\le257$. If the inequality holds for a prime $q\ge257$ and $r$ is the next
prime, then $r<2q$ by Bertrand's postulate and
$B_r=qB_q>9q^9>9(2q)^8>9r^8$. Thus every $n$ whose least nondividing prime is at least
$43$ has a representation.

For $q\le41$ put $M_q=\max_{\gcd(a,q)=1}P(a,q^3)$, the maximum over all reduced residue
classes modulo $q^3$. A sieve that retains the first prime in each reduced class gives
the exact values
\begin{center}
\begin{tabular}{r|rrrrrrr}
$q$ &2&3&5&7&11&13&17\\
$M_q$&17&109&1621&9371&47737&120401&349529\\ \hline
$q$&19&23&29&31&37&41&\\
$M_q$&540961&1494349&2367523&3812923&5897123&7802437&
\end{tabular}
\end{center}
An exception with least nondividing prime $q$ satisfies $B_q\le n\le M_q$. For
$q=23,29,31,37,41$ the lower endpoint exceeds the upper one, so it suffices to decide
representability for $1\le n\le540961$. For every prime $q$ with $q^3+2\le540961$ and
every residue $a$ modulo $q^3$, find the first prime $r$ in that class with
$r+q^3\le540961$, if one exists, and mark $r+q^3,r+2q^3,\dots$ up to the bound. The
marked values have explicit representations. Conversely every representation occurs
among the marks, since its prime is at least the first prime in the same class; this
step includes the class of $q$ itself. The $62$ unmarked values are those displayed.
Grouped by the least nondividing prime $q=2,3,5,7,11$ they number $6,11,25,15,5$, and
there are none for $q\ge13$.

The first program listed below carries out the sieve, the primorial checks and the
final enumeration; it tests each listed exception a second time by subtracting every
eligible positive multiple of each prime cube, and an independent bit-set sum of all
primes and all positive cube multiples agrees on every integer through $540961$. The
second program is an independent implementation, which tests the multiples of $B_q$ up
to $M_q$ coprime to $q$ one by one.
\end{proof}
\dirfiles{\artifact{grh-arithmetic-transfers/checks/prime_plus_cube_multiple.py},
\artifact{grh_cm_prime_applications/checks/cube.py}}

The largest exception can be checked by hand. For $n=10500$ the eligible primes $q$ are
$2,3,5,7,11,13,17,19$. The first four divide $n$, so the only candidate prime summand
is $q$ itself, and $q^3\nmid n-q$ in each case. For the other four an even multiplier
gives an even summand, and $n-2=10498$ is divisible by none of their cubes. For odd
multipliers the positive differences, grouped by $q$, are
\[
 \set{9169,6507,3845,1183},\quad\set{8303,3909},\quad\set{5587},\quad\set{3641},
\]
with prime divisors $53,3,5,7$, then $19,3$, then $37$ and $11$. On the other hand
$10501=5+1312\cdot2^3$.

\begin{corollary}[Prime squares; \stG]\label{dir:cor:square}
The positive integers that are not a prime plus a positive multiple of a prime square
are exactly $1,2,3,4,5,8,24$.
\end{corollary}

\begin{proof}
A multiple of $q^3$ is a multiple of $q^2$, so only the $62$ exceptions of
Theorem~\ref{dir:thm:cube} can fail. Direct subtraction of prime-square multiples on these
$62$ inputs leaves the seven displayed values.
\end{proof}

This is a second proof of the explicit conclusion of \cite{LeeOClarey2026Prime}, from
the same least-prime estimate.

\begin{theorem}[Prime fourth powers; \stG]\label{dir:thm:fourth}
Exactly $811$ positive integers fail to be a prime plus a positive multiple of a prime
fourth power. The largest is $1\,561\,560$.
\end{theorem}

\begin{proof}
Retain $q$ and $B_q$ from the proof of Theorem~\ref{dir:thm:cube}. For $q\ge53$,
Theorem~\ref{dir:thm:published}\ref{dir:it:leastprime} bounds the least prime congruent to $n$
modulo $q^4$ by $[4\varphi(q^4)\log q]^2<q^{10}$. Exact comparisons show $B_q>q^{10}$
for all primes from $53$ through $1031$, and Bertrand's postulate propagates the
inequality from any prime $q>2^{10}$ to its successor. Thus $n$ exceeds the least
congruent prime throughout this range.

For $q\le23$, enumerate all reduced classes modulo $q^4$ as before and test the
remaining multiples of $B_q$. The maxima of the least primes, in increasing order of
$q$, are
\[
 41,\ 647,\ 20063,\ 92369,\ 1189213,\ 2987429,\ 9951083,\ 19423919,\ 42026629,
\]
and the nine groups yield $12,48,230,255,212,46,7,1,0$ exceptions. When another prime
base $r$ dividing $n$ is tested, the only possible prime complement is $r$ itself,
which shortens the check.

For the six primes $29\le q\le47$ we have $4\log q<16$, so the least-prime bound covers
$n>[16q^3(q-1)]^2$. For each positive multiple $n$ of $B_q$ up to this bound with
$q\nmid n$, the checker exhibits a prime in the residue class of $n$ modulo $q^4$
strictly below $n$. The numbers of inputs in the six groups are
$516676,30586,4130,256,9,0$, and sieve lookup and trial division certify all the
primes. Independently, for every prime $r$ with $r^4\le2\,000\,000$, taking the first
prime $p$ in every occupied class modulo $r^4$ and marking $p+r^4,p+2r^4,\dots$ through
two million leaves the same $811$ unmarked integers.
\end{proof}
\dirfiles{\artifact{grh_cm_prime_applications/checks/fourth.py} (the list and the
computation), \artifact{grh_cm_prime_applications/checks/verify_exception_lists.py} (the
independent check)}

\begin{remark}\label{dir:rem:powers}
For each fixed exponent $j\ge2$ the same argument reduces $n=p+kq^j$ to finitely many
checks: the least-prime bound is at most $j^2q^{2j+2}$, while $B_q$ eventually exceeds
every fixed power of $q$.
\end{remark}

\subsection{Primes in geometric intervals and Pomerance's divisibility conjecture}
\label{dir:sec:geometric}
\sources{grh-arithmetic-transfers}

Barthel and M\"uller \cite{BarthelMuller2022Geometric} ask whether every interval
$[q^t,q^{t+1})$ meets every reduced residue class modulo $q$ in a prime, for integers
$q\ge2$ and $t\ge1$. Their Lemmas~5 and~6 cover all $q\le45000$ (with Bertrand's
postulate for $q=2$); their Remark~7 obtains all $t\ge3$ on Dirichlet GRH, and $t=2$
for $q\ge17386763$.

\begin{theorem}[Geometric intervals; \stG]\label{dir:thm:geometric}
For every $q\ge2$, every $a$ coprime to $q$ and every integer $t\ge2$, some prime
$p\equiv a\pmod q$ satisfies $q^t<p<q^{t+1}$. The same holds for $t=1$ whenever
$q\le50000$.
\end{theorem}

\begin{proof}
Dudek, Greni\'e and Molteni \cite[Theorem~1.1 and Table~1]{DudekGrenieMolteni2019Primes}
prove under Dirichlet GRH that for $q\ge3$ and $x\ge(18q\log q)^2$ there is a prime
$p\equiv a\pmod q$ within distance
\[
 H(x)=q\sqrt x\,\bigl(0.6265\log x+\log q+14\bigr)
\]
of $x$; here both their radius and their threshold have been enlarged using
$\varphi(q)\le q$. Their explicit formula is applied to primitive Dirichlet functions
and restores the factors at primes dividing the modulus, so the hypothesis is the one
supplied by \cref{src:thm:dgrh}.

For $q\ge50000$ write $L=\log q$ and $x_2=(18qL)^2$. The quotient $q/(378L^2+252L)$ is
increasing for $L>2$, by differentiation. At $q=50000$ we have $10<L<11$ and
$378L^2+252L<48510<50000$. Thus $x_2<q^3$ and
\[
 \frac{H(x_2)}{q^2}<18L(3L+14),\qquad q^2<x_2-H(x_2)<x_2+H(x_2)<q^3 .
\]
Indeed the upper endpoint divided by $q^2$ is at most $378L^2+252L$, and the lower one
exceeds $270L^2-252L>1$. For $t\ge2$ put $x_t=q^{t-2}x_2$. The function $H(x)/x$
decreases for $x\ge x_2$, again by differentiation. Therefore the interval of radius
$H(x_t)$ about $x_t$ is contained in the dilation by $q^{t-2}$ of the interval just
shown to lie inside $(q^2,q^3)$. This proves the statement for every $t\ge2$ and
$q\ge50000$.

For $45001\le q\le50000$ an exact sieve verifies $t=1$ and $t=2$ in every reduced
class (\artifact{grh-arithmetic-transfers/checks/progression_windows.cpp}). The
following numbers are the maxima of the \emph{first} prime witnesses when the search
starts at the indicated left endpoint, chosen so that all moduli share one sieve.
\begin{center}
\begin{tabular}{rrrr}
$t$ & search starts at & classes checked & largest witness\\ \hline
1 & $50001$ & $144390310$ & $9819037$\\
2 & $2500000001$ & $144390310$ & $2515292327$
\end{tabular}
\end{center}
For each row the program sieves $25\,000\,000$ consecutive integers by all primes up to
the square root of the last one, verifies that the segment lies inside the appropriate
geometric interval for each $q$, and traverses each reduced progression until it finds
a prime, using integer arithmetic only. For these moduli $t\ge3$ follows from Remark~7
of \cite{BarthelMuller2022Geometric}, and their results for $q\le45000$ complete the
proof. Neither endpoint of an interval can be a prime in the indicated class for
$q\ge2$.
\end{proof}

For a prime $p\ge3$ put
\[
 B(p)=\prod_{r<p,\ r\ \mathrm{prime}}r(r-1),\qquad
 D(p)=\frac{p-1}{\gcd(p-1,B(p))}.
\]
Pomerance \cite{Pomerance1974Carmichael} conjectured that $D(p)=1$ for every $p$, and
observed that if the least prime $P(1,q)$ in the class $1\bmod q$ satisfies
$q^2\mid P(1,q)-1$ for a prime $q$, then $p=P(1,q)$ violates the conjecture.

\begin{theorem}[The exact obstruction; \stG]\label{dir:thm:pomerance}
For every prime $p\ge3$, either $D(p)=1$, or $D(p)$ is a single prime $q$ for which
$p=P(1,q)$ and $v_q(p-1)=2$. In particular Pomerance's conjecture is equivalent to the
assertion that $q^2\nmid P(1,q)-1$ for every prime $q$.
\end{theorem}

\begin{proof}
Fix a prime $q\mid p-1$ and put $e=v_q(p-1)$. The factor $q$ itself occurs in $B(p)$.
If $q\le50000$, Theorem~\ref{dir:thm:geometric} for $t=1,\dots,e-1$ supplies $e-1$ distinct
primes $r<p$ with $q\mid r-1$; this handles all $e$, also when the list is empty.

Suppose $q>50000$. Above the class $1\bmod q$ there are two reduced classes modulo
$3q$, by the Chinese remainder theorem.
By Theorem~\ref{dir:thm:published}\ref{dir:it:leastprime} each of them contains a prime at
most $[2(q-1)\log(3q)]^2<q^3$; the last inequality follows from $4\log^2(3q)<q$, valid
at $50000$ and thereafter by monotonicity of $q/\log^2(3q)$. These two primes are
distinct. When $e\ge3$, use them first and then one prime from each of the $e-3$
intervals $[q^j,q^{j+1})$, $j=3,\dots,e-1$, in the class $1\bmod q$, from
Theorem~\ref{dir:thm:geometric}. They all lie below $q^e\le p-1$. Together with the factor $q$
they give $v_q(B(p))\ge e$.

The only possible shortage is at $e=2$. It occurs exactly when there is no earlier
prime $r\equiv1\pmod q$, that is, when $p=P(1,q)$, and the shortage is then exactly one
factor $q$. If $q<s$ were two deficient primes at $p$, we would have
$q^2s^2\le p-1<p=P(1,q)\le[(q-1)\log q]^2$ by the same least-prime bound, forcing
$s<\log q<q$, a contradiction.
\end{proof}

\begin{corollary}[Frequency of exceptions; \stG]\label{dir:cor:pomfreq}
The number of primes $p\le X$ with $D(p)\ne1$ is at most $\pi(\sqrt{(X-1)/2})$, hence
$O(\sqrt X/\log X)$; the exceptions have relative density zero among the primes.
\end{corollary}

\begin{proof}
By Theorem~\ref{dir:thm:pomerance} an exceptional $p$ is determined by its prime $q$. Both $p$
and $q$ are odd and $q^2\mid p-1$, so $2q^2\le p-1$.
\end{proof}

\begin{proposition}[A finite exclusion; \stU]\label{dir:prop:pomfinite}
For every prime $p\le2\cdot10^{14}+1$, $D(p)=1$. Moreover $P(1,q)<q^2$ for all $664579$
primes $q\le10^7$.
\end{proposition}

\begin{proof}
The program \artifact{grh-arithmetic-transfers/checks/least_primes.cpp} sieves the
prime moduli $q\le10^7$ and tests the candidates $kq+1$ in increasing order by exact
trial division. The largest least prime encountered is $1753933943$, and the largest
multiplier $k$ is $234$, first attained at $q=4722079$. For every prime $q\le10^7$ and
every $e\ge2$ with $q^e\le2\cdot10^{14}$, the same program verifies that the
$(e-1)$st prime in the class $1\bmod q$ lies below $q^e$; this checks $671350$ prime
powers. If $q$ is odd and $q^2\mid p-1$ then $2q^2\le p-1$, so these ranges include all
pairs $(q,e)$ relevant for $p\le2\cdot10^{14}+1$, also for $q=2$. The factor $q$ and the
$e-1$ earlier primes supply $q^e$ in $B(p)$.
\end{proof}

Consequently a counterexample to Pomerance's conjecture requires a prime $q>10^7$ and
$p=1+2kq^2\le[(q-1)\log q]^2$ for some positive integer $k$, with every smaller
candidate prime in the class $1\bmod q$ absent. Pomerance's assertion would rule out one
specific construction of a totient with a single preimage; it would not by itself
settle Carmichael's conjecture on the multiplicities of totients.

\subsection{The least prime-generating constants for five exponents}
\label{dir:sec:mills}
\sources{grh_cm_prime_applications}

For $\alpha>1$ call $A>1$ \emph{admissible} if $\floor{A^{\alpha^n}}$ is prime for
every integer $n\ge1$, and write $M_\alpha$ for the least admissible $A$ when one
exists. Chamberland and Straub \cite[Example~5.2]{ChamberlandStraub2026Prime} give the
digits $1.385510850556\ldots$ for $M_{9/4}$. Their construction produces an admissible
number, but an earlier prime, discarded by its selection rule, gives a strictly smaller
one.

\begin{theorem}[Certified digits; \stG]\label{dir:thm:mills}
The least admissible numbers for the following exponents exist and begin with the
digits shown. The last column gives the level of a prime from whose floor interval an
infinite prime sequence is proved to continue. Only the Riemann hypothesis for $\zeta$
is used.
\[
\begin{array}{c@{\qquad}l@{\quad}r}
\toprule
\alpha & M_\alpha \text{ begins with} & \text{level}\\
\midrule
9/4 &1.374277825758088669776655315788864285945977 &7\\
11/5 &1.40371466815791665549209159126764528659 &7\\
21/10 &1.440684278165675686775820360228 &9\\
41/20 &1.478399833482410327000644670 &10\\
201/100 &1.51462541826615615161458 &13\\
\bottomrule
\end{array}
\]
\end{theorem}

The proof occupies the rest of this subsection.

\paragraph{Where the smaller branch starts.}
At level $n$, the floor $p$ means that $A$ belongs to the half-open interval
\[
 I_n(p)=\bigl[p^{\alpha^{-n}},(p+1)^{\alpha^{-n}}\bigr).
\]
A common way to make $I_{n+1}(p')$ fit inside $I_n(p)$ is to require $p'\ge p^\alpha$
and $p'+1<(p+1)^\alpha$. For nonintegral $\alpha$ the intervals can also overlap when
$p'=\floor{p^\alpha}<p^\alpha$, and an overlap can contain a complete interval at a
later level. This happens at $\alpha=9/4$. Any admissible $A$ satisfies $A^\alpha\ge2$
and $A^{\alpha^2}\ge5$, since $2^\alpha$ lies between $4$ and $5$ and $4$ is composite.
The floor at the third level can be $37$, since $37<5^\alpha<38$. At the fourth level,
\[
 3455<5^{81/16}<3456<3457<3458<38^{9/4}.
\]
The integers $3455$ and $3456$ are composite and $3457$ is prime. Throughout
$3457\le A^{\alpha^4}<3458$ the earlier floors are $2,5,37$: the lower bound implies
$A^{\alpha^2}>5$, the upper bound gives $A^{\alpha^3}<38$, and
$2^\alpha<5$, $38<6^\alpha$, $6<3^\alpha$ give the remaining endpoints. Thus the whole
fourth-level interval fits inside the earlier overlap. An infinite continuation from
this interval gives $M_{9/4}<38^{\alpha^{-3}}<41^{\alpha^{-3}}$, while the branch in
\cite[Example~5.2]{ChamberlandStraub2026Prime} has third floor $41$.

\paragraph{The lower endpoint.}
Suppose $L$ is a lower bound for every admissible $A$. At level $n$ the next lower
bound is
\[
 L'=\max\set{L,\,p^{\alpha^{-n}}},\qquad
 p=\text{the least prime at least }\floor{L^{\alpha^n}} .
\]
This follows from monotonicity of the floor, with no assumption on the earlier floors
of a competitor. For an upper bound, start from the same finite choice of primes and
intersect their floor intervals. The certificate checks that the intersection contains
a whole terminal interval, that its upper endpoint lies below every earlier upper
endpoint, and that this terminal interval admits a continuation. Thus the lower bound
and the constructed admissible branch enclose the infimum of the admissible numbers.
The infimum is itself admissible: take admissible parameters tending to it from the
right; for each fixed $n$, right continuity makes the floors eventually equal to the
floor at the limit, and the initial lower bound $2^{1/\alpha}$ is strictly greater
than $1$.

\paragraph{The tail.}
Carneiro, Milinovich and Soundararajan
\cite[Theorem~5]{CarneiroMilinovichSoundararajan2017Fourier} prove on the Riemann
hypothesis that for real $y\ge4$ there is a prime in $[y,y+\tfrac{22}{25}\sqrt y\log y]$.
For $y=x^\alpha$ and $\alpha>2$, Taylor's theorem gives
$(x+1)^\alpha-1-x^\alpha>\alpha x^{\alpha-1}$ for $x\ge1$. Hence, when $x^\alpha\ge4$,
\[
 x^{(\alpha-2)/2}\ge\tfrac{22}{25}\log x
 \quad\Longrightarrow\quad
 \text{there is a prime } p \text{ with } x^\alpha\le p<(x+1)^\alpha-1 .
\]
This is also Corollary~3.2 of \cite{ChamberlandStraub2026Prime}; compare their
Lemma~5.1. For the five rows take $B=24,30,60,120,800$ respectively. If
$h=(\alpha-2)/2$, then $Bh\log10>1$ and $10^{Bh}>3B$ in each case. The function
$x^h/\log x$ increases past $10^B$, and $\log10<3$, so the implication applies for every
$x\ge10^B$. Each prime $p$ beyond this threshold therefore has a successor whose floor
interval lies inside its own, strictly at the upper end. Iteration gives an increasing
sequence of lower endpoints with a finite limit. Every old upper endpoint exceeds a
later upper endpoint, so this limit lies in every one of the half-open intervals and
is admissible.

\paragraph{Finite certificates.}
In the first seven levels the construction uses the least allowable prime at or above
the floor of the current lower bound, and each of the five cases ends this stage with a
complete seventh-level interval. If the seventh prime is below $10^B$, one subsequently
chooses any certified prime with a complete interval inside the previous one until the
threshold is passed. These later choices establish an upper bound; computing the first
seven floors of the minimum does not require that every later choice be minimal. The
lower and upper bounds for $M_\alpha$ are the two endpoints of the seventh-level
interval, the upper one strict: the lower endpoint is a global lower bound by
construction, and a continuation inside the interval proves the upper bound.
\[
\begin{array}{c|rrrrr|r}
\toprule
\alpha &p_1&p_2&p_3&p_4&p_5& \text{digits of the terminal prime}\\
\midrule
9/4&2&5&37&3457&91637597&41\\
11/5&2&5&37&2819&38923271&37\\
21/10&2&5&29&1213&2993027&126\\
41/20&2&5&29&997&1403869&223\\
201/100&2&5&29&877&823117&1577\\
\bottomrule
\end{array}
\]
The certificate contains every prime, floor and primality witness. If $L=b^{\alpha^{-i}}$ and $\alpha=m/d$, the verifier
checks $f^{d^{n-i}}\le b^{m^{n-i}}<(f+1)^{d^{n-i}}$ to establish a floor $f$. It proves
every skipped candidate composite by an integer divisor or a failed strong
probable-prime test. Each large prime $P$ comes with a factored divisor $F\mid P-1$,
$F^2>P$, and for each prime factor $q\mid F$ a witness $a$ with
$a^{P-1}\equiv1\pmod P$ and $\gcd(a^{(P-1)/q}-1,P)=1$. Thus for each prime divisor $r$
of $P$ the order of $a$ modulo $r$ contains the full $q$-part of $P-1$, so
$r\equiv1\pmod F$ for every $r\mid P$, whereas a composite $P$ would have a prime
divisor at most $\sqrt P$. The verifier recurses on $q$ and uses trial division below
$10^4$. Its decimal enclosures use repeated, outward-rounded rational roots, so the
table uses neither probabilistic primality tests nor floating-point comparisons. This
completes the proof of Theorem~\ref{dir:thm:mills}.
\dirfiles{\artifact{grh_cm_prime_applications/checks/mills_data.json} (certificate),
\artifact{grh_cm_prime_applications/checks/verify_mills.py} (verifier)}

\subsection{Toroczkai's prime-gap graphs}\label{dir:sec:graphs}
\sources{grh-exponent-six-and-transfers}

Let $p_n$ be the $n$th prime, put $p_0=1$ and $d_i=p_i-p_{i-1}$. A \emph{prime-gap
graph} on $n\ge2$ vertices is a simple graph on the vertex set $\set{1,\dots,n}$ with
$\deg(i)=d_i$. The degree sequence begins $1,1,2,2,4,2,4,2,4,6$, and the degree sum is
$p_n-1$. A matching is a set of edges with pairwise disjoint endpoints; $\nu(G)$ denotes
the largest size of a matching in $G$. Only the Riemann hypothesis for $\zeta$ is used
in this subsection.

\begin{theorem}[Toroczkai's two conjectures; \stG]\label{dir:thm:toroczkai}
For every $n\ge2$ a prime-gap graph on $n$ vertices exists, and every prime-gap graph
on $n$ vertices has a matching of size $d_{n+1}/2$.
\end{theorem}

\begin{proof}
Erd\H{o}s, Harcos, Kharel, Maga, Mezei and Toroczkai proved both assertions under the
Riemann hypothesis \cite[Theorems~1.1 and~1.6]{ErdosEtAl2024Sequence}, which
\cref{src:thm:dgrh} supplies. The second assertion implies the first: begin with a
single edge, and at each step replace $d_{n+1}/2$ disjoint edges by paths of length two
through a common new vertex. The old degrees stay fixed and the new vertex has the
required degree.
\end{proof}

Without the Riemann hypothesis, existence is known for $n\ge\exp(\exp(30.32))$ and the
matching assertion for $n\ge\exp(\exp(34.33))$ \cite{AggarwalFrotGouWang2025Explicit}.
The insertion step in the proof, which removes a matching of size $d_{n+1}/2$ and joins
its endpoints to the new vertex, is called an \emph{insertion} below.

Write $p=p_n$, $L=\log p$, and
\[
 N=\ceil{652L^2},\qquad B(p)=\ceil*{\frac{p-2}{4N}} .
\]

\begin{theorem}[Matching and block theorem; \stG]\label{dir:thm:block}
Every prime-gap graph $G$ on $n\ge2$ vertices has $\nu(G)\ge B(p)$. Whenever
$p_{n+k}-p\le2B(p)$, the next $k$ vertices can all be inserted using pairwise disjoint
sets of old endpoints: no two new vertices are adjacent, and every old vertex is
incident to at most one new edge. In particular, for each fixed $0<c<1/1304$ and all
sufficiently large $n$ this holds for $k=\floor{cp/L^3}$.
\end{theorem}

\begin{proof}
For an integer $T\ge3$ let $H(T)=\sum_{i\le n,\ d_i\ge T}d_i$. Theorem~1.8 of
\cite{ErdosEtAl2024Sequence} bounds the sum of the gaps of length at least $T$ that
start at a prime in $[x,2x]$ by $163x(\log x)^2/T$ under the Riemann hypothesis, for
every real $x\ge2$. Apply this at $x=p/2^j$ for $j\ge1$ while $x\ge2$. These intervals
cover every gap start up to $p$ except possibly $1,2,3$, whose following gaps have size
at most two. Since $\sum_{j\ge1}p/2^j=p$,
\begin{equation}\label{dir:eq:g1}
 H(T)\le163pL^2/T .
\end{equation}
Delete all vertices of degree at least $T$. This loses at most $H(T)$ edges and leaves
maximum degree at most $T-1$. Vizing's edge-colouring theorem partitions the remaining
edges into at most $T$ matchings. In particular
\begin{equation}\label{dir:eq:g2}
 \nu(G)\ge\ceil*{\frac{(p-1)/2-H(T)}{T}} .
\end{equation}
Taking $T=N$ gives $H(N)\le p/4$, hence $\nu(G)\ge B(p)$. The constant comes from the
leading lower bound $p/(2T)-163pL^2/T^2$, which is maximal at $T=652L^2$.

Partition a matching of size $(p_{n+k}-p)/2$ into pieces of sizes
$d_{n+1}/2,\dots,d_{n+k}/2$. For each piece, remove its edges and join their endpoints
to the corresponding new vertex. Distinct pieces have disjoint endpoints, so the
insertions can be made in any order. Finally $2B(p)\sim p/(1304L^2)$. For fixed
$c<c'<1/1304$ take $h=c'p/L^2$. Schoenfeld's estimate
$\pi(t)=\li(t)+O(\sqrt t\log t)$ under the Riemann hypothesis
\cite[Corollary~1]{Schoenfeld1976Sharper} gives
$\pi(p+h)-\pi(p)=(c'+o(1))p/L^3$ by subtraction at the two endpoints. Thus the first
$\floor{cp/L^3}$ new primes lie within the allowed distance once $p$ is large.
\end{proof}

\begin{proposition}[Restricted multiset extensions; \stU]\label{dir:prop:multiset}
Let $G$ be a finite simple graph and $e_1,\dots,e_k$ positive even integers. New
vertices of degrees $e_1,\dots,e_k$ can be appended to $G$ by insertions that use only
original edges with pairwise disjoint endpoints if and only if
$\sum_je_j\le2\nu(G)$.
\end{proposition}

\begin{proof}
Necessity holds because all removed edges form one matching. Sufficiency is the
partition construction in the proof of Theorem~\ref{dir:thm:block}.
\end{proof}

This is the restricted form of the extension question in Section~6 of
\cite{ErdosKharelMezeiToroczkai2023Degree}, which asks which multisets of prescribed
degrees can be appended to a fixed initial graph by degree-preserving steps. The
unrestricted problem permits reusing vertices and edges created by earlier steps; for
example, starting with a single edge one can add arbitrarily many vertices of degree
two by subdividing one current edge at each step. Theorem~\ref{dir:thm:block} gives a
quantitative sufficient condition for the unrestricted problem on every prime-gap
realization, with the stronger guarantee on which edges the extensions use.

\begin{example}[\stU]\label{dir:ex:graph}
For $n=10\,000$, $p=104729$ and $T=31$, the degree sum over the vertices of degree at
least $31$ is $9226$. Inequality \eqref{dir:eq:g2} gives
$\ceil{(52364-9226)/31}=1392$ disjoint edges in every realization, without an analytic
hypothesis. The next $236$ primes end at $107509$ and require
$(107509-104729)/2=1390$ edges in total; the following prime is $107563$. The program
\artifact{grh-exponent-six-and-transfers/checks/prime_gap_graphs.py} checks the degree
sums by sieving, constructs one realization by the Havel--Hakimi algorithm and carries
out all $236$ insertions, verifying simplicity, all $10236$ final degrees and
distinctness of the old endpoints
(\artifact{grh-exponent-six-and-transfers/data/graph-example.json}).
\end{example}

An infinite sequence of insertions can keep moving the same edge endpoint
indefinitely. A choice that stops moving old endpoints gives a connected countable
graph with the prescribed degrees, and one can control how long each vertex waits for
its final neighbours.

\begin{theorem}[Stabilization; \stG]\label{dir:thm:stabilize}
There is a sequence of connected prime-gap graphs $G_2,G_3,\dots$, each obtained from
the preceding one by an insertion, in which the neighbours of every fixed vertex
eventually become permanent. The limit is a connected, locally finite simple graph on
the positive integers with $\deg(i)=d_i$. Every limit edge $ij$ with $i<j$ satisfies
$j-i=O(j/\log j)$ as $j\to\infty$, uniformly in $i$.
\end{theorem}

\begin{proof}
Fix $A=100$ and a sufficiently large constant $P\ge e^A$. Use
Theorem~\ref{dir:thm:toroczkai} until the first stage with $p=p_n\ge P$. Connectivity persists
because each replaced edge has a new path between its original endpoints. At that
stage choose a spanning tree. Thereafter all of its edges, and the edges of the
successive trees defined below, are protected. Set
\[
 w(x)=\begin{cases}
 x,&0\le x<P,\\
 Ax/(A+\log(x/P)),&x\ge P,
 \end{cases}
 \qquad F(x)=x-w(x).
\]
At stage $n$ require that no edge used in the next insertion touch any of the first
$k=\pi(F(p))$ vertices. Both $w$ and $F$ are nondecreasing, and $F$ is $1$-Lipschitz.
For $p\ge P$, $w(p)\ge Ap/\log p$.

We show inductively that enough eligible edges remain. Let $u=p_k$ if $k>0$. At stage
$k$, all neighbours of the first $\pi(F(u))$ vertices had labels at most $k$ and were
permanently protected from that stage onward, by monotonicity of the cutoff. If $u<P$
this assertion is empty. Thus at least $(v-1)/2$ edges inside the first $k$ labels
still exist at stage $n$, where $v$ is the greatest prime at most $F(u)$, with $v=1$ if
there is none; this follows by summing the degrees of those earlier vertices and
allowing each internal edge to be counted twice. The degree sum over the first $k$
labels is $u-1$, so the number of edges incident to them is at most
$(u-1)-(v-1)/2$. Consequently at least $(p-2u+v)/2$ edges have both endpoints outside
the first $k$ labels.

By the prime-gap bound $d_{m+1}=O(\sqrt{p_m}\log p_m)$ under the Riemann hypothesis
\cite[Theorem~2.3]{ErdosEtAl2024Sequence}, $v\ge F(u)-O(\sqrt p\log p)$. As
$u\le F(p)$ and $F$ is $1$-Lipschitz, $-2u+F(u)\ge-2F(p)+F(F(p))$. Hence the number of
available edges, before protecting the tree or truncating high degrees, is at least
\[
 \frac{p-2F(p)+F(F(p))}{2}-O(\sqrt p\log p)
 =\frac{w(p)-w(F(p))}{2}-O(\sqrt p\log p).
\]
For $k=0$ the same inequality follows directly from the total of $(p-1)/2$ edges and
$0\le F(p)<2$. For $P\le x\le p$, differentiation gives $w'(x)=A(t-1)/t^2$ with
$t=A+\log(x/P)\ge A$, hence $w'(x)\ge w(p)/(2p)$; for $x<P$ the derivative equals one.
Integrating over $[F(p),p]$ gives
\[
 w(p)-w(F(p))\ge\frac{w(p)^2}{2p}.
\]
The part of the protected tree induced on the $n-k$ remaining vertices is a forest and
contains at most $n-k$ edges. For all sufficiently large $p$, uniformly in $P\le p$,
the prime-counting estimate under the Riemann hypothesis gives $n-k\le3w(p)/\log p$:
for $w(p)\ge p/2$ use $\pi(p)\sim p/\log p$; otherwise subtract the estimates for $\pi$
at $p$ and $p-w(p)\ge p/2$ and use $w(p)\ge Ap/\log p$ to absorb the error. Also
exclude the vertices of degree at least $T=\ceil{(\log p)^5}$. By \eqref{dir:eq:g1} the
resulting loss is at most $163p/(\log p)^3$ edges. The number of edges left is thus at
least
\[
 \frac{w(p)^2}{4p}-\frac{3w(p)}{\log p}-O(\sqrt p\log p)-\frac{163p}{(\log p)^3}
 \gg\frac{p}{(\log p)^2}.
\]
Here $w(p)^2/(4p)\ge25w(p)/\log p$. All implied constants are independent of $P$, so a
single sufficiently large $P$ makes these estimates valid throughout the induction.
Edge colouring now gives a matching of size $\gg p/(\log p)^7$, which exceeds
$d_{n+1}/2$. Make the insertion and add one edge incident to the new vertex to the
previously protected tree. The resulting trees are nested.

Since $F(p)\to\infty$, the neighbours of each fixed label eventually become permanent.
The union of the trees proves connectivity of the limit, and stabilization preserves
every finite degree. An edge $ij$ with $i<j$ is born at stage $j$, so for $p=p_{j-1}$
we have $i>\pi(F(p))$ once $p\ge P$. With $P$ fixed, $w(p)\sim100p/\log p$ and
$j-i\le1+\pi(p)-\pi(F(p))=O(p/(\log p)^2)=O(j/\log j)$.
\end{proof}

\begin{remark}[Which part uses the Riemann hypothesis]\label{dir:rem:graphrh}
A sufficiently late starting stage for Theorem~\ref{dir:thm:stabilize} can be treated with
unconditional estimates. Theorems~2.2 and~2.4 of \cite{ErdosEtAl2024Sequence} give
$d_{n+1}\le p^{5/8+o(1)}$ and $\sum_{i\le n}d_i^2\le p^{4/3+o(1)}$. With
$T=\ceil{p^{7/20}}$ one has
$H(T)\le T^{-1}\sum d_i^2\le p^{59/60+o(1)}=o(p/(\log p)^2)$. The gap estimate makes
the error from rounding $F(u)$ down to a prime $o(p/(\log p)^2)$ as well, and the
classical prime number theorem with remainder $O(x/(\log x)^4)$ suffices for
$n-k\le3w(p)/\log p$. After the same exclusions a matching has size
$\gg p^{13/20}/(\log p)^2$, exceeding $d_{n+1}/2$. Thus any sufficiently late connected
starting realization admits the construction by classical bounds alone. The start at
two vertices comes from Theorem~\ref{dir:thm:toroczkai}, and the uniform matching bound with
$T=\ceil{652L^2}$ uses the tail estimate \eqref{dir:eq:g1}.
\end{remark}

\subsection{Open questions}\label{dir:sec:open}
\sources{grh-transfer; grh-lean-transfers; grh-arithmetic-transfers; grh_cm_prime_applications; grh-exponent-six-and-transfers}

\begin{enumerate}
\item \emph{Additive cutoffs.} For each fixed $k\ge2$ the argument of
 Theorem~\ref{dir:thm:cube} reduces $n=p+mr^k$ to a finite search (Remark~\ref{dir:rem:powers}).
 What controls the growth with $k$ of the largest exception, which is $24$, $10500$
 and $1\,561\,560$ for $k=2,3,4$?
\item \emph{The first geometric interval.} Does every interval $[q,q^2)$ contain a
 prime in every reduced class modulo $q$ for $q>50000$? Is there a prime $q$ for which
 the least prime $1\bmod q$ is already $1\bmod q^2$? By Theorem~\ref{dir:thm:pomerance} the
 second question is equivalent to the failure of Pomerance's divisibility conjecture,
 and a positive answer to the first would rule out every such $q$.
\item \emph{Prime-generating constants near two.} The certificates of
 Theorem~\ref{dir:thm:mills} treat selected rational exponents. Can a common finite
 construction control a whole interval of exponents, including the changes where a
 powered endpoint crosses an integer? At exponent exactly two the factor
 $x^{(\alpha-2)/2}$ no longer grows, and the method does not decide whether the
 recursion taking the least prime above the square runs forever within the required
 intervals. Whether every interval between consecutive $\alpha$th powers contains a
 prime at $\alpha=2.01$, or even at $\alpha=2.2$, is a separate question.
\item \emph{Matchings in prime-gap graphs.} What is the asymptotic minimum of the
 largest matching among all prime-gap graphs on $n$ vertices? Theorem~\ref{dir:thm:block}
 gives a lower bound of order $p_n/(\log p_n)^2$; a better bound lengthens the block of
 insertions that use disjoint old endpoints.
\item \emph{Norm-Euclidean fields of degree eleven.} Which of the cyclic fields of
 degree $11$ and conductor $23$, $67$, $331$ are norm-Euclidean?
 Theorem~\ref{dir:thm:euclid} supplies no positive certificate.
\end{enumerate}
